\chapter{Carpal Tunnel Syndrome}

\section*{Definition and Overview}
Carpal tunnel syndrome is the most common entrapment neuropathy, caused by compression of the median nerve as it passes through the carpal tunnel -- a narrow channel at the front of the wrist formed by the wrist bones below and the tough flexor retinaculum ligament above. The median nerve carries sensation to the thumb, index, middle, and part of the ring finger, and supplies some of the small muscles of the thumb, so compression typically produces tingling, numbness, and, in later stages, weakness in these areas.

The condition usually develops gradually and may affect one or both hands. Symptoms are often mild and intermittent at first, frequently worse at night, and many people find shaking the hand or hanging it down brings relief. Left untreated, persistent compression can damage the nerve permanently, leading to loss of feeling and thumb weakness that may not fully recover.

\section*{Epidemiology}
Carpal tunnel syndrome is common, affecting roughly 3--6\% of adults at some point in life, and it is one of the most frequently reported work-related nerve complaints. It is about two to three times more common in women than in men and most often begins between the ages of 40 and 60.

Risk is higher in people who use their hands repeatedly or with forceful or vibratory movements, such as some assembly, manual, and computer-based work, although the strength of the occupational link is debated. Pregnant women and people with diabetes, obesity, hypothyroidism, rheumatoid arthritis, or kidney disease are also more likely to develop the condition. It can occur in one hand or both, and the dominant hand is more often affected.

\section*{Aetiology and Pathophysiology}
Carpal tunnel syndrome occurs when pressure rises inside the carpal tunnel and squeezes the median nerve. The tunnel is formed by the wrist bones at the back and the flexor retinaculum ligament at the front, and it also contains the tendons that bend the fingers. Anything that reduces the space in the tunnel or increases the bulk of its contents can raise pressure:

\begin{itemize}
    \item \textbf{Repetitive or forceful hand use:} sustained flexing and gripping can irritate the tendon sheaths and cause them to thicken
    \item \textbf{Fluid retention:} pregnancy, weight gain, and some medicines can swell the tissues
    \item \textbf{Inflammation:} rheumatoid arthritis and other inflammatory joint conditions
    \item \textbf{Underlying illness:} diabetes, hypothyroidism, obesity, and kidney disease increase susceptibility
    \item \textbf{Injury:} wrist fractures or dislocations that narrow the tunnel
    \item \textbf{Cysts or growths:} ganglion cysts occasionally press on the nerve
    \item \textbf{Idiopathic:} in most cases, no single cause is identified
\end{itemize}

Compression interferes with the blood supply to the nerve and damages the insulating myelin layer, slowing the nerve's signals. If compression persists, the nerve fibres themselves are damaged, causing weakness and wasting of the thumb muscles.

\section*{Clinical Features}
\begin{itemize}
    \item Tingling, pins and needles, or numbness in the thumb, index, middle, and ring finger, sparing the little finger
    \item Symptoms often worse at night, sometimes waking the person from sleep
    \item A feeling that the hand is "going to sleep", relieved by shaking or hanging the hand down
    \item Symptoms brought on by driving, holding a phone, reading, or other activities with the wrist bent
    \item Pain that may spread up the forearm and, occasionally, to the shoulder
    \item A dull ache or burning sensation in the hand and fingers
    \item Clumsiness and dropping objects in later stages
    \item Weakness and wasting of the fleshy mound at the base of the thumb (thenar eminence) in severe cases
    \item Reduced ability to distinguish fine touch or hot from cold in the affected fingers
\end{itemize}

\section*{Differential Diagnosis}
\begin{itemize}
    \item \textbf{Cervical radiculopathy:} nerve root compression in the neck, for example from a disc, which can mimic hand symptoms
    \item \textbf{Peripheral neuropathy:} generalised nerve disease, such as diabetic neuropathy, which usually affects both feet as well
    \item \textbf{Thoracic outlet syndrome:} compression of nerves or vessels between the neck and armpit
    \item \textbf{De Quervain's tenosynovitis:} tendon inflammation at the base of the thumb causing wrist pain
    \item \textbf{Arthritis of the wrist or base of the thumb}
    \item \textbf{Ulnar nerve compression} at the elbow or wrist, which affects the little finger and ring finger
    \item \textbf{Fibromyalgia and other widespread pain conditions}
\end{itemize}

\section*{When to Seek Medical Care}
See a doctor if tingling, numbness, or pain in the hand is persistent, worsening, wakes you repeatedly, or is affecting work and daily activities. Early treatment improves the chance of full recovery.

\textbf{Call 000 (or local emergency services) for:} sudden weakness of the hand, rapid loss of feeling, or a hand that becomes pale, cold, or blue. These may signal a different, more urgent problem such as a vascular or spinal emergency.

\section*{Diagnosis and Investigations}
\begin{itemize}
    \item \textbf{History and examination:} the distribution of symptoms and the effect on hand function are usually characteristic
    \item \textbf{Phalen's test:} holding the wrists in full flexion for about a minute reproduces symptoms in many affected people
    \item \textbf{Tinel's sign:} tapping over the wrist reproduces tingling in the fingers
    \item \textbf{Sensory testing:} checking for reduced feeling and two-point discrimination in the affected fingers
    \item \textbf{Nerve conduction studies and electromyography (EMG):} measure the speed of the median nerve and the health of the muscle it supplies; useful for confirming the diagnosis and grading severity
    \item \textbf{Ultrasound of the wrist:} can show a swollen median nerve or narrowing of the tunnel
    \item \textbf{Blood tests:} may be arranged to look for diabetes, thyroid disease, or inflammation when an underlying cause is suspected
\end{itemize}

\section*{Management}
\begin{itemize}
    \item \textbf{Wrist splinting:} wearing a splint at night (and sometimes during the day) keeps the wrist in a neutral position and relieves pressure; often the first-line treatment
    \item \textbf{Activity modification:} reducing repetitive, forceful, or sustained gripping and taking regular breaks
    \item \textbf{Cold packs and over-the-counter pain relief} for flare-ups
    \item \textbf{Corticosteroid injection:} an injection into the carpal tunnel can reduce inflammation and give weeks to months of relief; also useful as a diagnostic test and a bridging treatment
    \item \textbf{Treating underlying conditions:} controlling diabetes, treating hypothyroidism, or managing inflammatory arthritis
    \item \textbf{Physiotherapy and occupational therapy:} hand exercises, ergonomic advice, and workplace modifications
    \item \textbf{Surgery (carpal tunnel release):} dividing the ligament over the tunnel to relieve pressure; offered when conservative treatment fails or when weakness or nerve damage is present
    \item \textbf{Follow-up:} repeat nerve studies or review to check the response to treatment
\end{itemize}

\section*{Complications}
\begin{itemize}
    \item Permanent numbness, tingling, and loss of fine touch in the fingers
    \item Weakness and wasting of the thumb muscles, reducing grip strength and hand function
    \item Chronic pain and disability affecting work and daily life
    \item Clumsiness leading to falls, injuries, or dropping objects
    \item Treatment risks: corticosteroid injection carries a small risk of nerve damage or infection, and surgery rarely causes scar tenderness, nerve injury, or recurrent symptoms
\end{itemize}

\section*{Prognosis and Outcomes}
Many people with mild or intermittent symptoms improve with splinting and activity changes, and symptoms related to pregnancy often settle after the birth. Up to half of affected people improve without surgery, but about half of those with persistent symptoms eventually consider surgery. Carpal tunnel release is highly effective: the majority of people gain good relief of numbness and pain, although severe or long-standing cases may be left with some permanent loss of feeling or weakness. Recurrence after surgery is uncommon but possible.

\section*{Prevention and Self-Care}
\begin{itemize}
    \item Keep the wrist in a neutral (straight) position during activities and avoid sustained bending
    \item Take regular breaks from repetitive hand and wrist tasks
    \item Use ergonomic keyboards, mouse rests, and tools designed to reduce wrist strain
    \item Use the whole arm rather than just the wrist when working, and switch hands where possible
    \item Control diabetes, weight, and other conditions that raise risk
    \item If pregnant, mention hand symptoms to your doctor, since they often improve after delivery
    \item Manage early symptoms with a night splint and activity changes before they become established
    \item See a doctor early if symptoms persist, to avoid long-term nerve damage
\end{itemize}

\section*{Sources and Further Reading}
The following organisations provide reliable information on carpal tunnel syndrome, its diagnosis, and its treatment.

\begin{itemize}
    \item National Health Service. \textit{NHS conditions guide on carpal tunnel syndrome.} https://www.nhs.uk/conditions/
    \item Mayo Clinic. \textit{Carpal tunnel syndrome: symptoms, causes, and treatment.} https://www.mayoclinic.org/
    \item MSD Manual. \textit{Carpal tunnel syndrome -- professional overview.} https://www.msdmanuals.com/
    \item MedlinePlus. \textit{Consumer health information on carpal tunnel syndrome.} https://medlineplus.gov/
    \item National Institute for Health and Care Excellence. \textit{Clinical guidance on carpal tunnel syndrome.} https://www.nice.org.uk/
    \item Centers for Disease Control and Prevention. \textit{Workplace health information on ergonomics and hand disorders.} https://www.cdc.gov/
\end{itemize}
